% 系外卫星（综述）
% license CCBYSA3
% type Wiki

本文根据 CC-BY-SA 协议转载翻译自维基百科\href{https://en.wikipedia.org/wiki/Exomoon}{相关文章}。

\begin{figure}[ht]
\centering
\includegraphics[width=8cm]{./figures/f73ad0be1621979a.png}
\caption{候选系外卫星 Kepler-1625b I 环绕其母行星运行的艺术想象图 $^{[1]}$} \label{fig_Exomoo_1}
\end{figure}
系外卫星（exomoon），或称太阳系外卫星（extrasolar moon），是指环绕系外行星或其他非恒星的太阳系外天体运行的天然卫星。$^{[2]}$

受限于现有观测技术，系外卫星的探测与确认十分困难，$^{[3]}$ 截至目前尚无被正式确认的系外卫星发现。$^{[4]}$ 不过，诸如开普勒（Kepler）任务的观测已经发现了若干候选体。$^{[5][6]}$ 此外，通过微引力透镜方法，还探测到两个可能环绕流浪行星运行的系外卫星候选体。$^{[7][8]}$ 2019 年 9 月，天文学家报告称，对“塔比星”（Tabby’s Star）观测到的变暗现象，可能是由一颗孤立系外卫星被破坏后产生的碎片所致。$^{[9][10][11]}$ 一些系外卫星可能具备成为地外生命栖息地的潜力。$^{[2]}$ 在双流浪行星系统 2MASS J11193254−1137466 AB 的其中一个组分周围，还发现了一颗可能的系外卫星；如果该天体真实存在，则可能具备宜居条件。$^{[12]}$
\subsection{定义与命名}
传统用法中，“卫星”通常指环绕行星运行的天体，但近年来发现的棕矮星及其行星尺度的伴随天体，使行星与卫星之间的界限变得模糊，其根源在于棕矮星本身质量较低。对此，国际天文学联合会（IAU）通过如下声明加以澄清：“凡真实质量低于氘热核聚变临界质量、并环绕恒星、棕矮星或恒星遗骸运行，且其与中心天体的质量比低于 L4/L5 不稳定性阈值（$M/M_{\text{central}} < 2/(25+\sqrt{621})$）的天体，定义为行星。”$^{[13]}$

IAU 的该定义并未涉及那些质量低于棕矮星、且低于氘燃烧极限的自由漂浮天体（这类天体通常被称为自由漂浮行星、流浪行星、低质量棕矮星或孤立行星质量天体）的卫星命名规范。在相关文献中，这些天体的卫星通常仍被称为系外卫星。$^{[7][8][12]}$

系外卫星的命名规则通常沿用其母体天体的名称，并在其后加上一个大写罗马数字。例如，Kepler-1625b 环绕 Kepler-1625（亦称 Kepler-1625a）运行，而 Kepler-1625b 本身可能拥有一颗名为 Kepler-1625b I 的卫星（目前尚未发现 Kepler-1625b II，也尚不清楚 Kepler-1625b I 是否拥有次级卫星）。
\subsection{特征}
任何系外卫星的特征都可能各不相同，正如太阳系中的各类卫星一样。对于运行在其恒星宜居带内的系外巨行星而言，存在这样一种可能性：其所拥有的、尺寸接近类地行星的卫星，或许具备支持生命存在的条件。$^{[14][15]}$（需要进一步澄清）

2019 年 8 月，天文学家报告称，位于系外行星 WASP-49b 系统中的一颗系外卫星候选体，可能具有火山活动。$^{[16]}$
\subsection{轨道倾角}
对于由撞击形成、且距离其恒星不太远的类地行星卫星，如果行星—卫星之间的距离较大，那么在恒星潮汐力的作用下，卫星的轨道平面通常会趋向于与行星绕恒星运行的轨道平面保持一致；但如果行星—卫星距离较小，则其轨道可能具有明显的倾角。对于气态巨行星，其卫星的轨道往往与巨行星的赤道面保持一致，因为这些卫星形成于行星周围的环行星盘中。$^{[17]}$
\subsection{近恒星行星周围卫星的缺乏}
在近恒星、且具有近圆轨道的行星中，行星自转会逐渐减慢并最终发生潮汐锁定。随着行星自转减速，其同步轨道半径会向外移动。对于与恒星发生潮汐锁定的行星而言，卫星围绕该行星达到同步轨道的位置，往往位于行星希尔球之外。希尔球是行星引力相对于恒星占主导地位、从而能够束缚其卫星的空间区域。位于行星同步轨道半径以内的卫星，会因潮汐作用逐渐向内螺旋并最终坠入行星。因此，如果同步轨道位于希尔球之外，那么所有卫星最终都会螺旋坠入行星之中；如果同步轨道在三体问题意义下并不稳定，则卫星在到达同步轨道之前就会逃逸出行星轨道。$^{[17]}$

一项关于潮汐诱导迁移的研究为系外卫星稀少这一现象提供了一个可行解释。研究表明，母行星的物理演化（例如其内部结构与尺度）在决定卫星最终命运方面起着关键作用：同步轨道可能只是暂态状态，卫星可能被滞留在半渐近的半长轴位置，甚至被直接抛射出系统，而在此过程中还可能出现其他效应。这些因素反过来会对系外卫星的探测产生重大影响。$^{[18]}$
\subsection{探测方法}
理论上认为，许多系外行星周围都可能存在系外卫星。$^{[14]}$ 尽管行星搜寻者利用宿主恒星的多普勒光谱法已取得巨大成功，$^{[19]}$ 但这种技术并不能用于发现系外卫星。这是因为，行星及其附属卫星共同造成的恒星光谱位移，在观测上与一个绕宿主恒星运行的单一点质量天体并无区别。鉴于这一局限，人们提出了多种其他探测系外卫星的方法，包括：
\begin{itemize}
\item 直接成像
\item 微引力透镜
\item 脉冲星计时
\item 凌日时间变化效应
\item 凌日法
\end{itemize}
\subsubsection{直接成像}
由于恒星与系外行星之间亮度差异极大，以及行星本身体积小、辐射弱，对系外行星进行直接成像极具挑战性。在大多数情况下，这一问题对系外卫星而言更为严峻。然而，有理论指出，受到强烈潮汐加热的系外卫星，其亮度可能与某些系外行星相当。潮汐力会通过差异引力耗散能量，从而加热卫星。木星的卫星木卫一（Io）便是典型例子，其火山活动正是由潮汐加热所驱动。如果一颗系外卫星受到足够强的潮汐加热，且与其恒星的距离足够远，使得卫星自身的光不至于被恒星淹没，那么诸如詹姆斯·韦布空间望远镜这样的设备，理论上就有可能对其进行成像。$^{[20]}$
\subsubsection{宿主行星的多普勒光谱法}
多普勒光谱法是一种间接探测方法，通过测量绕行天体引起的速度变化以及相应的光谱位移来发现行星。$^{[21]}$ 该方法也被称为径向速度法，对主序星尤为有效。在若干案例中，人们已成功部分获取了系外行星的光谱，例如 HD 189733 b 和 HD 209458 b。然而，与恒星光谱相比，行星光谱更容易受到噪声影响，因此其光谱分辨率以及可提取的谱线数量，都远低于执行系外行星\subsubsection{多普勒光谱分析所需的水平。
来自宿主行星磁层的射电辐射}
在运行过程中，木卫一的电离层与木星磁层相互作用，产生摩擦电流，从而引发射电辐射，这种辐射被称为“由木卫一控制的分米波辐射”。研究人员认为，在已知系外行星附近寻找类似的射电辐射，可能成为预测其是否存在其他卫星的关键线索。$^{[22]}$
\subsubsection{微引力透镜}
2002 年，Cheongho Han 与 Wonyong Han 提出可利用微引力透镜效应来探测系外卫星。$^{[23]}$ 他们发现，即便透镜事件涉及角半径较小的源恒星，由于严重的有限源效应，卫星信号在透镜光变曲线中会被显著抹平，因此探测难度极大。
\subsubsection{脉冲星计时}
2008 年，澳大利亚莫纳什大学的 Lewis、Sackett 和 Mardling 提出使用脉冲星计时方法来探测脉冲星行星的卫星。$^{[24]}$ 他们将该方法应用于 PSR B1620-26 b 的情形，发现如果该行星拥有一颗稳定的卫星，且卫星与行星的距离约为行星绕脉冲星轨道半径的五十分之一，同时其质量不低于行星质量的 5\%，则这种卫星是可以被探测到的。
\subsubsection{凌日时间效应}
2007 年，物理学家 A. Simon、K. Szatmáry 和 Gy. M. Szabó 发表了一篇研究简报，题为《通过光度凌日时间变化确定“系外卫星”的大小、质量和密度》。$^{[25]}$

2009 年，David Kipping 发表论文 $^{[3][26]}$，系统阐述了如何将多次观测到的凌日中点时间变化（TTV：当行星—卫星系统的质心相对于视线方向处于近乎垂直状态时，由于行星在质心前后摆动所引起）与凌日持续时间变化（TDV：当卫星—行星连线大致沿视线方向时，由于行星沿凌日路径相对于质心运动所引起）结合起来，从而产生一种独特的系外卫星信号。此外，该研究还表明，可以利用这两种效应同时确定系外卫星的质量以及其绕行星运行的轨道距离。

在随后的一项研究中，Kipping 得出结论：利用 TTV 与 TDV 效应，开普勒空间望远镜有能力探测到位于宜居带内的系外卫星。$^{[27]}$
\subsubsection{凌日法（恒星—行星—卫星系统）}
当一颗系外行星从宿主恒星前方经过时，来自恒星的观测光度会出现一次微小的下降。凌日法是目前探测系外行星最成功、最灵敏的方法。这一效应（也称掩星）与行星半径的平方成正比。如果一颗行星及其卫星同时从宿主恒星前方经过，那么二者都应在观测光曲线中造成亮度下降。$^{[28]}$
在凌日过程中还可能发生行星—卫星相互掩食的现象，$^{[29]}$ 但此类事件本身发生概率较低。
\subsubsection{凌日法（行星—卫星系统）}
如果宿主行星能够被直接成像，那么系外卫星的凌日也有可能被观测到。当系外卫星从宿主行星前方经过时，来自该行星的光度会出现一次微小下降。$^{[29]}$
对于被直接成像的系外行星以及自由漂浮行星，其系外卫星被凌日探测到的概率和发生率预计都较高。利用这一方法，詹姆斯·韦布空间望远镜理论上可以探测到小至木卫一（Io）或土卫六（Titan）规模的卫星，但该方法需要大量观测时间。$^{[12]}$
\subsubsection{轨道采样效应}
如果将玻璃瓶对着光线观察，会发现瓶身中央比边缘更容易透光。类似地，在对一颗卫星绕行星运动进行位置采样时，采样点在轨道两端会比在中间更为密集。如果一颗卫星绕行一颗发生恒星凌日的行星运行，那么该卫星也会发生恒星凌日；当观测数据量足够多时，这种在轨道两端的“堆积效应”可能会在凌日光曲线中显现出来。恒星越大，为了检测到这种堆积效应所需的观测次数就越多。开普勒望远镜的数据可能足以利用轨道采样效应探测红矮星周围的系外卫星，但对类太阳恒星而言数据量仍然不足。$^{[30][31]}$
\subsubsection{围绕热木星的间接探测}

类比木星的卫星木卫一，有理论认为，围绕热木星运行并受到强烈潮汐加热的系外卫星，其表面可能覆盖着富含钠或钾的特殊包层。这种包层可通过地面光谱仪进行探测，从而间接指示这些系外卫星的存在。$^{[32][33]}$
\subsubsection{围绕白矮星的间接探测}
白矮星的大气中有时会被金属元素“污染”，并且在少数情况下，白矮星周围还存在碎屑盘。通常认为，这种污染源自小行星或彗星，但过去也有人提出，遭到潮汐撕裂的系外卫星同样可能是白矮星污染的来源之一。$^{[34]}$
2021 年，Klein 及其合作者发现白矮星 GD 378 和 GALEXJ2339 的铍元素污染程度异常之高。研究人员认为，为了产生如此过量的铍，氧、碳或氮原子必须经历了与质子能量达 MeV 量级的碰撞。$^{[35]}$
在一种解释情景中，这种铍的过量来自被潮汐撕裂的系外卫星：在该模型中，一颗巨行星绕白矮星运行，其周围存在一个形成卫星的冰质盘。巨行星强大的磁场会加速恒星风中的质子，并将其引导进入该盘。被加速的质子与盘中的水冰发生碰撞，通过散裂反应生成铍、硼和锂等元素。由于这些元素在恒星核聚变过程中会被破坏，因此在宇宙中相对稀少。若卫星微体在这种盘中形成，其铍、硼和锂的丰度将明显偏高。该研究还预测，例如土星的一些中等大小卫星（如土卫一 Mimas）也应富集 Be、B 和 Li。$^{[36]}$
